\documentclass[11pt]{article}
\usepackage[margin=1in]{geometry}
\usepackage{amsmath,amssymb,amsthm,hyperref}
\newtheorem{theorem}{Theorem}
\newtheorem{lemma}[theorem]{Lemma}
\newtheorem{corollary}[theorem]{Corollary}
\newcommand{\E}{\mathbb E}
\newcommand{\Pp}{\mathbb P}
\newcommand{\R}{\mathbb R}
\title{A one-neuron target that defeats wide sign-only Bayes}
\date{18 November 2026; independent research draft}
\emergencystretch=3em
\author{ChatGPT-6 Astra\thanks{Independent AI research. Sole author. Verification scope and provenance are stated in the accompanying research collection.}}
\hypersetup{pdftitle={A one-neuron target that defeats wide sign-only Bayes},pdfauthor={ChatGPT-6 Astra}}
\begin{document}
\maketitle

\section{Statement and interpretation}
This note gives a particularly transparent separation between sign-only
Bayesian conditioning and fixed positive training margins. It uses a
cosine activation, not the repository's usual tanh or ReLU activation.
The connection between Gaussian kernels and parity hardness is classical;
the precise online construction and its finite-neural-width transfer below
still need a targeted novelty check. The mathematical claims should be
assessed on their own, without an assertion that the phenomenon is new.

Let $X=\{-1,1\}^d$ and $f(x)=\prod_{i=1}^d x_i$. Consider
\[
 F_N(x)=\frac1N\sum_{j=1}^N a_j\cos(w_j\cdot x+b_j),
\]
with independent $a_j\sim N(0,1)$, $w_j\sim N(0,\sigma^2I_d)$,
and $b_j$ uniform modulo $2\pi$. The teacher has just one neuron:
\[
 f(x)=\cos\left(\frac\pi2\sum_i x_i-\frac{d\pi}{2}\right).
\]
Its weights are bounded by $\pi/2$ and its phase can be reduced modulo
$2\pi$. Thus its description and number of parameters are $O(d)$.

The sign-only posterior conditions on $f(x)F_N(x)>0$ at observed inputs.
There is no initial global output restriction. Present the cube in
binary lexicographic order by writing $x_i=1-2z_i$, $z_i\in\{0,1\}$.

\begin{theorem}[Wide sign-only failure]
For every $d\ge1$ and $\sigma>0$, all sufficiently large finite widths
$N$ have posterior-majority prediction errors on every input after the
first in this order: exactly $2^d-1$ forced errors, irrespective of how
the first tie is resolved. Fresh posterior sampling has expected errors
at least $0.49(2^d-1)$ for all sufficiently large $N$. When $\sigma$ is
fixed, an explicit width threshold of the form $\exp(O(d))$ suffices
for either statement, with constants depending on $\sigma$.
\end{theorem}

\begin{corollary}[Width and margin separations]
At width one, the same sign-only prior has an $O(d\log(d+1))$ total
expected sampling-mistake bound, uniform over all input orders, and a
bound of the same order for any fixed-confidence mistakes. Thus widening
can change polynomial online behavior into exponential behavior.

If instead training requires $f(x)F_N(x)>1/2$, the quantitative entropy
projection argument in the companion note gives polynomial
total mistake bounds uniform over every $N\ge1$. The same prior and
activation are used. This latter statement uses mean-field normalization
$1/N$ and a fixed, nonvanishing training margin.
\end{corollary}

The sign-only rule is invariant under positive rescaling of $F_N$.
It therefore sees a central-limit Gaussian process at large width.
A fixed positive mean-field margin is a large-deviation constraint and
does not remain in that Gaussian-process regime. This is the mechanism
behind the separation, rather than an alteration of the target.

There is also a necessary width scale. At any width, keep one teacher
neuron in the width-one box below and constrain all other readouts to
$|a_j|\le1/(8N)$, leaving their incoming weights unrestricted. Their
total unnormalized interference is less than $1/8$, so every target sign
is still correct. This event has negative log prior mass at most
$O_\sigma(d\log(d+1)+N\log(N+1))$. If all $M-1$ later posterior-majority
predictions are wrong, each next correct-label probability is at most
$1/2$, so the full-label evidence is at most $2^{-(M-1)}$. Comparing
the two estimates forces $N=\Omega_\sigma(2^d/d)$ for large $d$.
Thus the width scale for the all-but-first failure is exponential in
dimension, up to the different exponential constants in the upper bound.

\section{An exact Gaussian-process calculation}
Put $r=e^{-2\sigma^2}\in(0,1)$ and $M=2^d$. The normalized covariance
of the cosine features is
\[
 K(z,z')=r^{|z\mathbin\triangle z'|}.
\]
Indeed averaging over the phase removes the sum-frequency term, and
the Gaussian characteristic function gives
$2\E\cos(w\cdot x+b)\cos(w\cdot x'+b)=K(z,z')$.
Let $G$ be the centered Gaussian vector with covariance $K$, and let
$H_z=(-1)^{|z|}G_z$ be its signed version. Its covariance is
\[
 \widetilde K=
 \begin{pmatrix}1&-r\\-r&1\end{pmatrix}^{\otimes d}.
\]
Write $q=\sqrt{1-r^2}$. In lexicographic order the Cholesky factor is
$\bigl(\begin{smallmatrix}1&0\\-r&q\end{smallmatrix}\bigr)^{\otimes d}$.
Its inverse is
$\bigl(\begin{smallmatrix}1&0\\r/q&1/q\end{smallmatrix}\bigr)^{\otimes d}$.
Consequently the exact autoregressive representation is
\begin{equation}\label{eq:ar}
 H_z=-\sum_{y\subsetneq z}r^{|z\setminus y|}H_y+q^{|z|}\varepsilon_z,
 \qquad \varepsilon_z\stackrel{\rm iid}{\sim}N(0,1).
\end{equation}
Every proper subset $y$ precedes $z$; the innovation is independent of
the entire preceding Gaussian vector, not merely the displayed subsets.
If all previous signed outputs are positive and $z\ne\varnothing$,
the conditional mean in \eqref{eq:ar} is strictly negative. Averaging
over the truncated distribution of preceding values therefore gives
\[
 \Pp(H_z>0\mid H_y>0\text{ for all preceding }y)<1/2.
\]
This proves the $M-1$ majority errors for the Gaussian process. The
conditioning probabilities are all positive because $K\succ0$.
The multivariate central limit theorem already transfers each strict
inequality to all sufficiently large neural widths; there are only
$M-1$ inequalities. The next section supplies quantitative bounds.

\section{A finite-width transfer using only scalar concentration}
Set
\[
 a=(1-r)^d,\qquad b=(1+r)^d.
\]
Every principal covariance matrix of $\widetilde K$ has eigenvalues
between $a$ and $b$. A lower bound for the probability of the full
positive orthant, and hence every prefix orthant, is
\begin{equation}\label{eq:orthantlower}
 p_*:=\left(\frac{a}{2\pi e b}\right)^{M/2}.
\end{equation}
To see this, integrate its Gaussian density on $[0,\sqrt a]^M$:
the determinant is at most $b^M$ and the quadratic form at most $M$.

Let $h_j(x)=\sqrt2\cos(w_j\cdot x+b_j)$ and
$K_N=N^{-1}\sum_j h_jh_j^T$. Conditional on all incoming weights and
phases, $\sqrt{2N}F_N$ is exactly Gaussian with covariance $K_N$.
Since every product $h_j(x)h_j(y)$ lies in $[-2,2]$, scalar Hoeffding
and a union bound give
\begin{equation}\label{eq:covconcentration}
 \Pp\{\|K_N-K\|_{\rm op}>\epsilon a\}
 \le2M^2\exp\left(-\frac{N\epsilon^2a^2}{8M^2}\right).
\end{equation}
Here entrywise error $\epsilon a/M$ implies the indicated operator
bound. On its complement,
$(1-\epsilon)K\preceq K_N\preceq(1+\epsilon)K$.

For any cone $C\subseteq\R^M$, not necessarily convex, these Loewner
bounds imply
\begin{equation}\label{eq:coneratio}
 e^{-2M\epsilon}\Pp_K(C)\le\Pp_{K_N}(C)
                 \le e^{2M\epsilon}\Pp_K(C),\qquad 0<\epsilon\le1/2.
\end{equation}
Indeed compare densities with those of the scalar multiples
$(1\pm\epsilon)K$, whose cone probabilities equal $\Pp_K(C)$;
the determinant factors are bounded by
$[(1+\epsilon)/(1-\epsilon)]^{M/2}\le e^{2M\epsilon}$.
Prefix training events and their intersections with a query sign are
cones. If the bad covariance probability is at most $q_*$, the neural
posterior correct-sign probability is therefore at most
\begin{equation}\label{eq:posteriortransfer}
 \frac{e^{4M\epsilon}}{1-q_*}\,p_{\mathrm{GP}}
 +\frac{q_*e^{2M\epsilon}}{(1-q_*)p_*},
\end{equation}
where $p_{\mathrm{GP}}$ is the Gaussian-process correct-sign probability.
This explicitly accounts for reweighting the incoming features by the
rare training event; convergence of unconditioned covariances alone
would not suffice.

For the sampling assertion use $\epsilon=1/(1000M)$ and
$q_*=p_*/1000$. Substituting $p_{\mathrm{GP}}\le1/2$ in
\eqref{eq:posteriortransfer} makes the result less than $0.51$.
It suffices to take
\[
 N\ge\frac{8M^2}{\epsilon^2a^2}
                    \log\frac{2M^2}{q_*}.
\]
At fixed $r$, its logarithm is $O(d)$. Fresh draws along this deterministic
order are independent, so ordinary Bernoulli concentration also gives
an exponential-in-$M$ high-probability lower bound on the number of
mistakes, if desired.

\section{A uniform strict majority gap}
The following elementary bound avoids leaving the majority width
threshold unquantified. For every nonempty query,
\begin{equation}\label{eq:gap}
 p_{\mathrm{GP}}\le\frac12-\delta_d,
 \qquad\delta_d=\frac{r^d a^2}{1000b^{3/2}\sqrt M}.
\end{equation}
Let $V$ be the Gaussian vector at the $k\le M$ preceding inputs,
already signed by parity, conditioned on its positive orthant. Whitening
turns this event into a cone. The Gaussian radius is independent of its
direction and is unchanged by conditioning on a cone. Consequently
$\E[\|V\|^2\mid V>0]\le bk$. With conditional probability at least
$3/4$, all coordinates other than the first have norm at most
$2\sqrt{bM}$.

Conditional on these other coordinates, the first coordinate is a normal
$N(\mu,s^2)$ truncated to $(0,\infty)$. The inverse covariance bounds
give $s^2\in[a,b]$ and
$\mu\ge-T$, where $T=2b^{3/2}\sqrt M/a$, on that event. Set
$u=a/[4(T+\sqrt b)]$. The standard-normal hazard bound
$\phi(t)/\overline\Phi(t)\le t_++1$ implies
\[
 \frac{\overline\Phi((u-\mu)/s)}{\overline\Phi(-\mu/s)}
 \ge\exp\left(-\frac{uT}{a}-\frac{u^2}{2a}-\frac{u}{\sqrt a}\right)
 \ge e^{-1}.
\]
Thus $\Pp(V_0\ge u\mid V>0)\ge3/(4e)$. In \eqref{eq:ar}, the
coefficient of $V_0$ is at least $r^d$, and the innovation variance is
at most one. Hence the excess wrong-sign probability over $1/2$ is
at least
\[
 \frac{3}{4e}[\Phi(r^du)-1/2]
 \ge\frac{3\phi(1)}{4e}r^du
 \ge\frac{\phi(1)}{16e}
           \frac{r^da^2}{b^{3/2}\sqrt M}>\delta_d.
\]
We used $u\le1$, $T+\sqrt b\le3b^{3/2}\sqrt M/a$, and
$\phi(1)/(16e)>1/1000$. The hazard bound follows from the usual Mills
ratio lower bound $\overline\Phi(t)/\phi(t)\ge1/(t+1)$ for $t\ge0$;
for $t<0$ the hazard is at most $\sqrt{2/\pi}<1$.

Now take $\epsilon=\delta_d/(64M)$ and $q_*=\delta_dp_*/16$ in
\eqref{eq:covconcentration}--\eqref{eq:posteriortransfer}. The latter
is less than $1/2-\delta_d/2$. This proves all $M-1$ majority mistakes
at an explicit finite width threshold whose logarithm is again $O(d)$
for fixed $r$. The constants are large and no sharp transition width
is claimed.

\section{The small-width and positive-margin guarantees}
At width one, take the incoming weight center $w_i^*=\pi/2$, phase
$b^*=-d\pi/2$ modulo $2\pi$, and readout in $[1,2]$. If
$|w_i-w_i^*|\le1/(8d)$ and the circular phase error is at most $1/8$,
then the phase error at every input is at most $1/4$. Thus
$f(x)F_1(x)\ge\cos(1/4)>1/2$ everywhere. This box has prior probability
at least $\exp[-O_\sigma(d\log(d+1))]$. The ordinary evidence-potential
argument therefore bounds total expected sign mistakes, and fixed-
confidence mistakes, by $O_\sigma(d\log(d+1))$ on every sequence.

For the positive-margin large-width theorem, an explicit mean teacher
law is especially simple. Take $w\sim N(w^*,I_d/d)$, phase uniformly
within circular distance $1/4$ of $b^*$, and independently
$a\sim N(\mu,1)$, where
\[
 \mu=\frac{2e^{1/2}}{\operatorname{sinc}(1/4)}<4,
 \qquad\operatorname{sinc}(t)=\sin(t)/t.
\]
Its mean output is exactly $2f(x)$ everywhere. Its entropy relative to
the stated prior is
\[
 K=\frac d2\left[\frac{(\pi/2)^2}{\sigma^2}
       +\frac1{d\sigma^2}-1+\log(d\sigma^2)\right]
       +\log(4\pi)+\frac{\mu^2}{2}.
\]
This is $O_\sigma(d\log(d+1))$. The bounded cosine feature gives all
the same output-moment and conditional-Gaussian-readout estimates as
the tanh theorem. Its covariance $K/2$ is positive definite, proving
full feature span. Thus the sufficient width and mistake bound are
polynomial. For smaller widths pad and amplify the single-neuron teacher
as in the companion note; the product-box prior cost is polynomial in
$N,d$, and the split at the polynomial sufficient width yields a
uniform polynomial over every $N\ge1$.

\section{A related separation for ordinary Gaussian tanh}
The cosine example gives the stronger all-but-first majority conclusion,
but a superpolynomial confidence and sampling separation for ordinary
tanh already follows by combining an existing repository lower bound
with this round's new positive-margin upper bound. We state this
consequence explicitly so the role of the activation is not misleading.

Use $F_N=N^{-1}\sum_j a_j\tanh(u_j\cdot(1,x))$, with
$a_j\sim N(0,\tau^2)$ and $u_j\sim N(0,\tau^2I_{d+1}/d)$ at fixed
$\tau>0$. Let $f$ be full parity and $d\ge25$. Put
\[
 T_d=1+\left\lfloor\frac1{6000}
           \binom d{\lfloor\sqrt d/5\rfloor}\right\rfloor.
\]
The repository paper \emph{Parity under standard-scaling neural-network
Gaussian processes}, Theorem ``Sign-only learner'' and its network
corollary, prove the following existing fact: after any history of at
most $T_d-1$ inputs, at least a quarter of the cube has limiting
correct-sign posterior probability at most $13/20$.
Its constant $C_0<1000$ makes the displayed conservative threshold valid.
Positive rescaling of the output leaves sign-only conditioning unchanged,
so its Gaussian-process limit applies to the mean-field normalization
used here as well. The first-layer gain is chosen to match the same
incoming variance; the cited bound is uniform in this gain.

Recursively choose an input in this bad quarter using the limiting
posterior. It is unseen, since observed labels have posterior probability
one. This yields one fixed sequence of $T_d$ distinct inputs, depending
on $d$ and the prior but not on sufficiently large $N$. Posterior
convergence at its finitely many histories makes the finite-width
correct probabilities less than $2/3$ at every round once $N$ is large
enough. Consequently sign-only learning has $T_d$ confidence failures
and at least $T_d/3$ expected fresh-sampling mistakes. Here
$T_d=\exp(\Omega(\sqrt d\log d))$; no finite-width threshold is claimed.
This is an online consequence of the repository's existing theorem,
not a new proof of its parity lower bound, and it does not imply a
majority-error lower bound.

For comparison, a small integer-weight tanh teacher is explicit. With
$S=\sum_i x_i$ and $c_d=(1+(-1)^d)/2$, set
\[
 F_T(x)=c_d\tanh(d)+\sum_{j=1}^d(-1)^{d-j}
                 \tanh\bigl(d(S+d-2j+1)\bigr).
\]
Replacing tanh by sign makes the expression exactly parity. Each
nonconstant argument has absolute value at least $d$, so its uniform
error is at most $2(d+1)e^{-2d}<1/2$. This is a width-at-most-$d+1$
teacher with integer parameter magnitudes at most $d^2$ and positive
margin greater than $1/2$. The companion entropy-projection theorem
therefore gives polynomial total mistakes over every width $N\ge d+1$
when the same prior instead conditions only on observed margins
$f(x)F_N(x)>1/2$. Its explicit constants are polynomial in the teacher
weight bound, so using $B_0=d^2$ preserves polynomial dependence on $d$.

Thus changing zero training margin to a fixed positive one separates
superpolynomial and polynomial lifetime sampling/confidence bounds even
for the ordinary Gaussian tanh prior. Neither side imposes a global
margin restriction at unobserved inputs.

\section{Verification status and remaining questions}
Independent review and formal verification remain outstanding. The
analogous majority-error lower bound for ordinary tanh or ReLU, and a
practical transition width, remain open.
\end{document}
